% =====================================================================
\chapter{Proposiciones, conectivas y tablas de verdad}
\label{tema:1}
% =====================================================================

Una demostración es una cadena de afirmaciones enlazadas por reglas de razonamiento. Antes de estudiar las técnicas de demostración conviene fijar qué es una afirmación, cómo se combinan varias mediante conectivas y cuándo dos combinaciones dicen lo mismo. Este tema desarrolla la lógica proposicional con el enfoque del capítulo~2 de~\cite{hammack} (§2.1–2.6 y §2.11) y del capítulo~1 de~\cite{velleman}; la parte algebraica sigue la §3.4 de~\cite{mcs}.

\section{Proposiciones y conectivas}

\begin{definicion}[Proposición]\label{def:proposicion}
Una \emph{proposición} es un enunciado del que tiene sentido afirmar que es verdadero o que es falso, y que es exactamente una de las dos cosas. Su \emph{valor de verdad} es $\V$ (verdadero) o $\F$ (falso).
\end{definicion}

\begin{ejemplo}\label{ej:proposiciones}
«$7$ es un número primo» es una proposición verdadera; «$2+2=5$» es una proposición falsa; «todo número par mayor que~2 es suma de dos primos» (conjetura de Goldbach) es una proposición, aunque hoy no se sepa cuál es su valor de verdad. En cambio, «¿qué hora es?», «cierre la puerta» y «$x+1=3$» no son proposiciones: las dos primeras no afirman nada y la tercera depende del valor de $x$, que no está fijado. Enunciados como el último se llaman \emph{predicados} y se estudian en el \cref{tema:2}.
\end{ejemplo}

A partir de proposiciones $P$ y $Q$ se forman proposiciones compuestas mediante \emph{conectivas lógicas}. El valor de verdad de la proposición compuesta depende solo de los de $P$ y $Q$, y se describe con una \emph{tabla de verdad}.

\begin{definicion}[Negación, conjunción y disyunción]\label{def:conectivas}
Sean $P$ y $Q$ proposiciones.
\begin{enumerate}
  \item La \emph{negación} $\neg P$ («no $P$») es verdadera si y solo si $P$ es falsa.
  \item La \emph{conjunción} $P\wedge Q$ («$P$ y $Q$») es verdadera si y solo si $P$ y $Q$ son ambas verdaderas.
  \item La \emph{disyunción} $P\vee Q$ («$P$ o $Q$») es verdadera si y solo si al menos una de las dos es verdadera.
\end{enumerate}
\end{definicion}

\begin{table}[h]
\centering
\caption{Tablas de verdad de las conectivas básicas.}\label{tab:conectivas}
\begin{tabular}{@{}cc|ccccc@{}}
\toprule
$P$ & $Q$ & $\neg P$ & $P\wedge Q$ & $P\vee Q$ & $P\Rightarrow Q$ & $P\Leftrightarrow Q$ \\
\midrule
$\V$ & $\V$ & $\F$ & $\V$ & $\V$ & $\V$ & $\V$ \\
$\V$ & $\F$ & $\F$ & $\F$ & $\V$ & $\F$ & $\F$ \\
$\F$ & $\V$ & $\V$ & $\F$ & $\V$ & $\V$ & $\F$ \\
$\F$ & $\F$ & $\V$ & $\F$ & $\F$ & $\V$ & $\V$ \\
\bottomrule
\end{tabular}
\end{table}

La disyunción matemática es \emph{inclusiva}: $P\vee Q$ es verdadera también cuando lo son las dos. La disyunción \emph{exclusiva}, «o $P$ o $Q$, pero no ambas», se escribe $P\xor Q$ y se estudia en el \cref{eje:1-xor}.

\section{Condicional y bicondicional}

\begin{definicion}[Condicional]\label{def:condicional}
El \emph{condicional} $P\Rightarrow Q$ («si $P$, entonces $Q$») es falso únicamente cuando $P$ es verdadera y $Q$ es falsa. $P$ se llama \emph{hipótesis} o antecedente, y $Q$, \emph{tesis} o consecuente.
\end{definicion}

La tabla del condicional (\cref{tab:conectivas}) se entiende como una promesa: «si apruebas el examen, te invito a cenar» solo se incumple si se aprueba y no hay cena. Si no se aprueba, la promesa no se ha roto, haya cena o no. Por eso un condicional con hipótesis falsa es verdadero; se dice que es \emph{vacuamente} verdadero.

El condicional $P\Rightarrow Q$ admite varias lecturas equivalentes, que conviene reconocer porque aparecen constantemente en los enunciados de los teoremas~\cite[§2.3]{hammack}:
\begin{center}
\begin{tabular}{@{}l@{}}
«si $P$, entonces $Q$»; \quad «$Q$ si $P$»; \quad «$P$ solo si $Q$»; \\
«$P$ es condición suficiente para $Q$»; \quad «$Q$ es condición necesaria para $P$».
\end{tabular}
\end{center}

\begin{definicion}[Bicondicional]\label{def:bicondicional}
El \emph{bicondicional} $P\Leftrightarrow Q$ («$P$ si y solo si $Q$») es verdadero cuando $P$ y $Q$ tienen el mismo valor de verdad. Se lee también «$P$ es condición necesaria y suficiente para $Q$».
\end{definicion}

\begin{definicion}[Recíproco, contrarrecíproco e inverso]\label{def:reciproco}
Dado el condicional $P\Rightarrow Q$, se llaman \emph{recíproco} a $Q\Rightarrow P$, \emph{contrarrecíproco} a $\neg Q\Rightarrow\neg P$ e \emph{inverso} a $\neg P\Rightarrow\neg Q$.
\end{definicion}

\begin{ejemplo}\label{ej:reciproco}
Sea $n$ un entero y considérese «si $n$ es múltiplo de~4, entonces $n$ es par», que es verdadero. Su recíproco, «si $n$ es par, entonces es múltiplo de~4», es falso: $n=2$ lo refuta. Su contrarrecíproco, «si $n$ no es par, entonces no es múltiplo de~4», es verdadero. La \cref{prop:leyes}(v) muestra que un condicional y su contrarrecíproco siempre tienen el mismo valor de verdad, mientras que el recíproco puede no tenerlo.
\end{ejemplo}

\section{Fórmulas, tautologías y contradicciones}

\begin{definicion}[Fórmula proposicional]
Las \emph{fórmulas proposicionales} sobre las variables $P_1,\dots,P_n$ son las expresiones que se obtienen a partir de las variables aplicando un número finito de veces las conectivas $\neg$, $\wedge$, $\vee$, $\Rightarrow$ y $\Leftrightarrow$, con paréntesis para indicar el orden. Una \emph{asignación} da a cada variable un valor de verdad, y la tabla de verdad de la fórmula lista su valor para cada asignación.
\end{definicion}

\begin{proposicion}\label{prop:filas}
Una fórmula con $n$ variables tiene una tabla de verdad de $2^{n}$ filas.
\end{proposicion}

\begin{proof}
Una asignación es una elección independiente de $\V$ o $\F$ para cada una de las $n$ variables. Por el principio del producto, hay $2\cdot 2\cdots 2=2^{n}$ elecciones.
\end{proof}

\begin{definicion}[Tautología y contradicción]
Una fórmula es una \emph{tautología} si es verdadera para todas las asignaciones, y una \emph{contradicción} si es falsa para todas.
\end{definicion}

\begin{ejemplo}\label{ej:tercero-excluido}
$P\vee\neg P$ (principio del tercero excluido) es una tautología, y $P\wedge\neg P$ es una contradicción: en cada fila exactamente una de $P$ y $\neg P$ es verdadera. Más interesante es la fórmula $\bigl((P\Rightarrow Q)\wedge P\bigr)\Rightarrow Q$. Su única forma de ser falsa sería con $Q$ falsa y $(P\Rightarrow Q)\wedge P$ verdadera; pero entonces $P$ es verdadera y $P\Rightarrow Q$ también, lo que obliga a que $Q$ sea verdadera. Por tanto es una tautología; es la regla de \emph{modus ponens} de la \cref{sec:inferencia}.
\end{ejemplo}

El razonamiento del ejemplo anterior evita escribir la tabla completa: para comprobar que un condicional es tautología basta ver que la fila «hipótesis verdadera y tesis falsa» es imposible. Esta forma de pensar es ya una demostración por reducción al absurdo (\cref{tema:4}).

\section{Equivalencia lógica}

\begin{definicion}[Equivalencia lógica]\label{def:equivalencia}
Dos fórmulas $\Phi$ y $\Psi$ en las mismas variables son \emph{lógicamente equivalentes}, y se escribe $\Phi\eqlog\Psi$, si tienen el mismo valor de verdad en todas las asignaciones; equivalentemente, si $\Phi\Leftrightarrow\Psi$ es una tautología.
\end{definicion}

\begin{proposicion}[Leyes lógicas]\label{prop:leyes}
Para proposiciones cualesquiera $P$, $Q$ y $R$:
\begin{enumerate}
  \item doble negación: $\neg\neg P\eqlog P$;
  \item leyes de De Morgan: $\neg(P\wedge Q)\eqlog\neg P\vee\neg Q$ \ y \ $\neg(P\vee Q)\eqlog\neg P\wedge\neg Q$;
  \item distributivas: $P\wedge(Q\vee R)\eqlog(P\wedge Q)\vee(P\wedge R)$ \ y \ $P\vee(Q\wedge R)\eqlog(P\vee Q)\wedge(P\vee R)$;
  \item condicional como disyunción: $P\Rightarrow Q\eqlog\neg P\vee Q$;
  \item contrarrecíproco: $P\Rightarrow Q\eqlog\neg Q\Rightarrow\neg P$;
  \item negación del condicional: $\neg(P\Rightarrow Q)\eqlog P\wedge\neg Q$;
  \item bicondicional: $P\Leftrightarrow Q\eqlog(P\Rightarrow Q)\wedge(Q\Rightarrow P)$.
\end{enumerate}
Además, $\wedge$ y $\vee$ son conmutativas y asociativas.
\end{proposicion}

\begin{proof}
Cada apartado se comprueba con una tabla de verdad. Se detallan (ii), (iv) y (v); los demás son análogos (\cref{eje:1-leyes}).

(ii) $\neg(P\wedge Q)$ es verdadera exactamente cuando $P\wedge Q$ es falsa, es decir, cuando al menos una de $P$ o $Q$ es falsa, es decir, cuando al menos una de $\neg P$ o $\neg Q$ es verdadera, que es la condición para que $\neg P\vee\neg Q$ sea verdadera. La segunda ley se razona igual intercambiando los papeles de «todas» y «al menos una».

(iv) $\neg P\vee Q$ es falsa exactamente cuando $\neg P$ y $Q$ son falsas, es decir, cuando $P$ es verdadera y $Q$ falsa. Es la misma única fila en la que $P\Rightarrow Q$ es falsa.

(v) Por (iv), $\neg Q\Rightarrow\neg P\eqlog\neg\neg Q\vee\neg P$, que por (i) y la conmutatividad equivale a $\neg P\vee Q$, y de nuevo por (iv), a $P\Rightarrow Q$.
\end{proof}

Las leyes de la \cref{prop:leyes} permiten transformar fórmulas sin construir tablas, igual que las propiedades de la suma y el producto permiten manipular expresiones algebraicas~\cite[§3.4]{mcs}.

\begin{ejemplo}\label{ej:simplificar}
Se simplifica $\neg(P\Rightarrow\neg Q)$:
\[
\neg(P\Rightarrow\neg Q)\ \eqlog\ P\wedge\neg\neg Q\ \eqlog\ P\wedge Q,
\]
por (vi) y (i). Así, negar «si llueve, no salgo» equivale a afirmar «llueve y salgo».
\end{ejemplo}

\section{Razonamientos válidos}\label{sec:inferencia}

\begin{definicion}[Razonamiento válido]
Un razonamiento con premisas $\Phi_1,\dots,\Phi_k$ y conclusión $\Psi$ es \emph{válido} si $(\Phi_1\wedge\dots\wedge\Phi_k)\Rightarrow\Psi$ es una tautología, es decir, si en toda asignación que hace verdaderas todas las premisas también es verdadera la conclusión.
\end{definicion}

\begin{proposicion}[Reglas de inferencia]\label{prop:inferencia}
Los siguientes razonamientos son válidos~\cite[§2.11]{hammack}:
\begin{center}
\begin{tabular}{@{}lll@{}}
\toprule
Nombre & Premisas & Conclusión \\
\midrule
\emph{Modus ponens} & $P\Rightarrow Q$, \ $P$ & $Q$ \\
\emph{Modus tollens} & $P\Rightarrow Q$, \ $\neg Q$ & $\neg P$ \\
Silogismo hipotético & $P\Rightarrow Q$, \ $Q\Rightarrow R$ & $P\Rightarrow R$ \\
Silogismo disyuntivo & $P\vee Q$, \ $\neg P$ & $Q$ \\
\bottomrule
\end{tabular}
\end{center}
\end{proposicion}

\begin{proof}
\emph{Modus ponens} se probó en el \cref{ej:tercero-excluido}. Para \emph{modus tollens}, si $P\Rightarrow Q$ y $\neg Q$ son verdaderas, por la \cref{prop:leyes}(v) también lo es $\neg Q\Rightarrow\neg P$, y \emph{modus ponens} da $\neg P$. Para el silogismo disyuntivo, si $\neg P$ es verdadera, $P$ es falsa, y como $P\vee Q$ es verdadera, ha de serlo $Q$. El silogismo hipotético es el \cref{eje:1-silogismo}.
\end{proof}

\begin{observacion}[Errores frecuentes]\label{obs:1-errores}
(a) Las \emph{falacias} más comunes imitan reglas válidas: de $P\Rightarrow Q$ y $Q$ \emph{no} se sigue $P$ (afirmar el consecuente), y de $P\Rightarrow Q$ y $\neg P$ \emph{no} se sigue $\neg Q$ (negar el antecedente). En ambos casos la asignación $P=\F$, $Q=\V$ hace verdaderas las premisas y falsa la conclusión.
(b) Demostrar el recíproco de un teorema no demuestra el teorema (\cref{ej:reciproco}).
(c) El «o» matemático es inclusivo; si se quiere excluir el caso «ambos», debe decirse explícitamente.
(d) $P\Rightarrow Q$ no afirma que $P$ sea verdadera ni que $P$ cause $Q$; solo excluye el caso $P$ verdadera y $Q$ falsa.
\end{observacion}

\begin{observacion}[Conexión con el curso]
Los enunciados de los módulos posteriores distinguen con cuidado entre condiciones necesarias y suficientes. Por ejemplo, en el módulo~4 el gradiente nulo es una condición \emph{necesaria} para que un punto interior sea mínimo de una función diferenciable, pero no suficiente; y el criterio de la hessiana definida positiva es \emph{suficiente}, pero no necesario. Confundir una condición con su recíproca es el error de la \cref{obs:1-errores}(b). En el módulo~13, la máscara causal de la atención es literalmente una tabla de valores de verdad: la posición $i$ puede atender a la $j$ si y solo si $j\le i$.
\end{observacion}

\begin{lecturas}
\item Hammack~\cite{hammack}, cap.~2 «Logic», §2.1–2.6 y §2.11 «Logical Inference».
\item Velleman~\cite{velleman}, §1.1 «Deductive Reasoning and Logical Connectives», §1.2 «Truth Tables» y §1.5 «The Conditional and Biconditional Connectives».
\item Lehman, Leighton y Meyer~\cite{mcs}, §3.1 «Propositions from Propositions», §3.3 «Equivalence and Validity» y §3.4 «The Algebra of Propositions».
\item Stanford CS103~\cite{cs103}, clase~3 «Propositional Logic».
\end{lecturas}

\section*{Ejercicios del tema 1}
\addcontentsline{toc}{section}{Ejercicios del tema 1}

\begin{ejercicio}{1}{eje:1-proposiciones}
Indique cuáles de los siguientes enunciados son proposiciones y, en su caso, su valor de verdad: (a) «$2^{10}>1000$»; (b) «$x^{2}\ge 0$»; (c) «el número $91$ es primo»; (d) «calcule $\sqrt{2}$ con tres decimales»; (e) «existe un número real $x$ con $x^{2}=-1$»; (f) «todo entero es par o impar».
\end{ejercicio}

\begin{ejercicio}{1}{eje:1-traducir}
Sean $P$: «$n$ es par», $Q$: «$n$ es múltiplo de~3» y $R$: «$n$ es múltiplo de~6», donde $n$ es un entero fijo. Escriba en símbolos: (a) «$n$ es múltiplo de~6 si y solo si es par y múltiplo de~3»; (b) «$n$ es par solo si es múltiplo de~6»; (c) «ser par no es suficiente para ser múltiplo de~6»; (d) «$n$ no es par ni múltiplo de~3». Indique cuáles de (a)–(c) son verdaderas para todo entero $n$.
\end{ejercicio}

\begin{ejercicio}{1}{eje:1-reciproco}
Para el enunciado «si una función $f\colon\R\to\R$ es derivable en $a$, entonces es continua en $a$», escriba el recíproco, el contrarrecíproco y el inverso, e indique cuáles son verdaderos. Para los falsos, dé una función que lo muestre.
\fuente{adaptado de Hammack~\cite{hammack}, §2.8 «More on Conditional Statements».}
\end{ejercicio}

\begin{ejercicio}{1}{eje:1-necesaria}
Para un entero $n$, clasifique cada condición como necesaria, suficiente, ambas o ninguna para que «$n$ sea múltiplo de~4»: (a) $n$ es par; (b) $n$ es múltiplo de~8; (c) $n^{2}$ es múltiplo de~16; (d) $n$ termina en~4.
\end{ejercicio}

\begin{ejercicio}{2}{eje:1-silogismo}
Construya la tabla de verdad de $\bigl((P\Rightarrow Q)\wedge(Q\Rightarrow R)\bigr)\Rightarrow(P\Rightarrow R)$ y concluya que el silogismo hipotético de la \cref{prop:inferencia} es válido.
\fuente{adaptado de Hammack~\cite{hammack}, §2.5 «Truth Tables for Statements».}
\end{ejercicio}

\begin{ejercicio}{2}{eje:1-xor}
La disyunción exclusiva $P\xor Q$ es verdadera si y solo si exactamente una de $P$ y $Q$ lo es. (a) Escriba su tabla de verdad. (b) Compruebe que $P\xor Q\eqlog(P\vee Q)\wedge\neg(P\wedge Q)$. (c) Compruebe que $P\xor Q\eqlog\neg(P\Leftrightarrow Q)$.
\end{ejercicio}

\begin{ejercicio}{2}{eje:1-equivalencias}
Decida con tablas de verdad si son equivalentes: (a) $P\Rightarrow(Q\Rightarrow R)$ y $(P\wedge Q)\Rightarrow R$; (b) $(P\Rightarrow Q)\Rightarrow R$ y $P\Rightarrow(Q\Rightarrow R)$; (c) $(P\vee Q)\Rightarrow R$ y $(P\Rightarrow R)\wedge(Q\Rightarrow R)$. Cuando no lo sean, indique una asignación que las distinga.
\fuente{adaptado de Velleman~\cite{velleman}, §1.5 «The Conditional and Biconditional Connectives».}
\end{ejercicio}

\begin{ejercicio}{2}{eje:1-algebra}
Usando solo las leyes de la \cref{prop:leyes} (sin tablas), demuestre que $\neg\bigl(P\vee(\neg P\wedge Q)\bigr)\eqlog\neg P\wedge\neg Q$ y que $(P\Rightarrow Q)\wedge(P\Rightarrow\neg Q)\eqlog\neg P$.
\fuente{adaptado de Lehman, Leighton y Meyer~\cite{mcs}, §3.4 «The Algebra of Propositions».}
\end{ejercicio}

\begin{ejercicio}{2}{eje:1-validez}
Decida si los razonamientos siguientes son válidos e identifique la regla o la falacia correspondiente. (a) «Si el modelo está sobreajustado, el error de validación sube. El error de validación sube. Luego el modelo está sobreajustado.» (b) «Si $n$ es primo y mayor que~2, entonces es impar. $n$ es par. Luego $n$ no es un primo mayor que~2.» (c) «El fallo está en los datos o en el código. No está en los datos. Luego está en el código.» (d) «Si $x>2$, entonces $x^{2}>4$. $x\le 2$. Luego $x^{2}\le 4$.»
\end{ejercicio}

\begin{ejercicio}{2}{eje:1-contar}
(a) ¿Cuántas tablas de verdad distintas puede tener una fórmula con $n$ variables? Calcule el número para $n=1,2,3$. (b) Para $n=2$, ¿cuántas de esas funciones son tautologías y cuántas contradicciones? (c) Escriba una fórmula para cada una de las cuatro funciones de una variable.
\end{ejercicio}

\begin{ejercicio}{2}{eje:1-leyes}
Complete la demostración de la \cref{prop:leyes} comprobando con tablas de verdad las leyes distributivas (iii) y las leyes (vi) y (vii).
\end{ejercicio}

\begin{ejercicio}{3}{eje:1-demorgan-n}
Sea $n\ge 2$ y sean $P_1,\dots,P_n$ proposiciones. Demuestre las leyes de De Morgan generalizadas
\[
\neg(P_1\wedge\dots\wedge P_n)\eqlog\neg P_1\vee\dots\vee\neg P_n,\qquad
\neg(P_1\vee\dots\vee P_n)\eqlog\neg P_1\wedge\dots\wedge\neg P_n,
\]
razonando directamente sobre los valores de verdad.
\end{ejercicio}

\begin{ejercicio}{3}{eje:1-completitud}
Demuestre que toda tabla de verdad con $n\ge 1$ variables es la de alguna fórmula que solo usa las conectivas $\neg$ y $\wedge$. (Indicación: construya primero una fórmula con $\neg$, $\wedge$ y $\vee$ que sea una disyunción de conjunciones, una por cada fila en la que la tabla vale $\V$, y elimine después $\vee$ con De Morgan.)
\fuente{adaptado de Lehman, Leighton y Meyer~\cite{mcs}, §3.4 «The Algebra of Propositions» (formas normales).}
\end{ejercicio}

\begin{ejercicio}{3}{eje:1-bicond-asoc}
Demuestre que el bicondicional es asociativo, $(P\Leftrightarrow Q)\Leftrightarrow R\eqlog P\Leftrightarrow(Q\Leftrightarrow R)$, probando que ambos miembros son verdaderos exactamente cuando el número de proposiciones falsas entre $P$, $Q$ y $R$ es par. Demuestre después, con una asignación concreta, que el condicional no es asociativo.
\end{ejercicio}
